\documentclass{tufte-handout}

\usepackage{lmodern}
\usepackage{amsmath}
\usepackage{amssymb}

\newcounter{example}
\newcounter{exc:first:proof}
\newenvironment{example}
{\refstepcounter{example}\begin{quote}
\textbf{Example \arabic{example}}}
{

$\square$\end{quote}}

\newenvironment{exercise}
{\refstepcounter{example}\begin{quote}
\textbf{Exercise \arabic{example}}}
{\end{quote}}

\newcommand{\RED}[1]{\textcolor{red}{#1}}
\newcommand{\TODO}[1]{\RED{TODO: #1}}

\newcounter{theExerciseCounter}

\title[You can Prove it!]{You can Prove it!}

\author{Elias Castegren \& Ellen Arvidsson}

\synctex=1

\begin{document}

\section{Suggested Solutions}

Note that these are \emph{suggested} solutions. There may be other
equally valid ways of writing these proofs.

\subsection{``What do you Mean by That?''}

\begin{enumerate}
\item All three members of TLC share the opinion that anyone who
  is a scrub is a guy who cannot get any love from them. We
  introduce a set $\mathit{TLC}$ of members of TLC, the predicates
  $\mathit{guy}$ and $\mathit{scrub}$, and a function
  $\mathit{love}$ such that $\mathit{love}(a, b)$ returns a
  natural number representing the amount of love $a$ is willing to
  give to $b$.
%
  We use the ``for all'' quantifier to introduce names $p$ and $q$
  for different people:
  \[
    \forall p. p \in \mathit{TLC} \implies
    (\forall q. scrub(q) \implies guy(q) \land love(p, q) = 0)
  \]
  Note that we can put both quantifiers in the beginning and turn
  the proposition into a single implication:
  \[
    \forall p, q. p \in \mathit{TLC} \land scrub(q) \implies
    guy(q) \land love(p, q) = 0
  \]
\item This exercise leaves a lot open for modeling. Here is one
  suggestion. We are assuming a pitch which is a one-dimensional
  line with player and ball positions as distinct points on this
  line. Let us put coordinates on the line so that we can talk
  about relative positions.
%
  We will put the center of the pitch on the $0$ coordinate, and
  the two goals on some coordinates $\mathit{goal}$ and
  $-\mathit{goal}$. We can sketch it as follows:

\begin{verbatim}
  <---|-------------|-------------|--->
    -goal           0            goal
\end{verbatim}

  For simplicity, let us define offside position for players
  trying to score goals on the positive side of the
  pitch\footnote{We can easily flip things around to get the
    definition for the other team}. We need to capture three
  properties for the position of some player $p$ and the ball. We
  will imagine a function $\mathit{pos}$ defined for both players
  and the ball:
  \begin{enumerate}
  \item Being in the opponents' half of the pitch
  \item Being closer to the opponents' goal than the ball is
  \item Being closer to the opponents' goal than the
    second-closest opponent
  \end{enumerate}
  The first property can be simply $\mathit{pos}(p) > 0$.
%
  The second property is also quite straightforward:
  $\mathit{pos}(p) > \mathit{pos}(\mathit{ball})$.
%
  The third property requires some more thought. For example, we
  could define the set of opposing players in front of player $p$
  and use the size of that set to check if there are fewer than
  two opponents in front. If we let $\mathit{opponents}(p)$ be the
  set of players on the opposing team, we can filter the players
  in front with the following set comprehension:
  $\{q~|~q\in\mathit{opponents}(p) ~\land~ pos(p)<pos(q)\}$. In
  the end we get the following solution:
\[
  \begin{array}{l}
    \forall p. \mathit{offside}(p) \iff\\
    \quad \mathit{pos}(p) > 0 ~\land\\
    \quad \mathit{pos}(p) > \mathit{pos}(\mathit{ball}) ~\land\\
    \quad \left|\{q~|~q\in\mathit{opponents}(p) ~\land~ pos(p)<pos(q)\}\right| < 2
  \end{array}
\]

% A different choice is to let $0$ be the coordinate of one of the
% goals and $\mathit{goal}$ be the coordinate of the other. This
% puts the center of the pitch at $\frac{\mathit{goal}}{2}$:

% \begin{verbatim}
%   <---|-------------|-------------|--->
%       0           goal/2         goal
% \end{verbatim}


Note that this is \emph{one} way of defining offside position.
Another choice of model may lead to another definition. In the
end, your definitions and level of abstraction depends on what you
want to do with the formalisation in the end.
\end{enumerate}

\subsection{``Oh Yeah? Prove it!''}

\begin{enumerate}
\setcounter{enumi}{2}
\item We want to show
  $\forall A~B~C. ~A \land (B \lor C) \iff (A \land B) \lor (A
  \land C)$. We start by assuming we have some propositions $A$,
  $B$ and $C$ and proceed to show
  $A \land (B \lor C) \iff (A \land B) \lor (A \land C)$. We
  proceed by cases over the two directions:
  \begin{itemize}
  \item[\textbf{Case 1:}] $A \land (B \lor C) \implies (A \land B) \lor (A \land C)$\\
    %
    We start by assuming $A \land (B \lor C)$. We need to
    show $(A \land B) \lor (A \land C)$.
    %
    From our assumption we get $A$ and $B \lor C$. We proceed by
    cases over $B \lor C$:
    \begin{enumerate}
    \item If $B$ holds then we can show that $A \land B$ holds
      since both $A$ and $B$ hold.
    \item If $C$ holds then we can show that $A \land C$ holds
      since both $A$ and $C$ hold.
    \end{enumerate}

  \item[\textbf{Case 2:}]
    $(A \land B) \lor (A \land C) \implies A \land (B \lor C)$\\
    %
    We start by assuming $(A \land B) \lor (A \land C)$. We need
    to show $A \land (B \lor C)$. We proceed by cases over
    $(A \land B) \lor (A \land C)$:
    \begin{enumerate}
    \item If $(A \land B)$ holds, then $A$ and $B$ hold. Now $A$
      holds by assumption and from $B$ we can show $B \lor C$.
    \item If $(A \land C)$ holds, then $A$ and $C$ hold. Now $A$
      holds by assumption and from $C$ we can show $B \lor C$.
    \end{enumerate}
    $\square$
  \end{itemize}

\item We want to show that $n^3 + 11n$ is divisible by $6$ for any
  natural number $n$.

  \textbf{Proof}\\
  The proof is by induction over $n$:
  \begin{itemize}
  \item \textbf{Base Case}: $n = 0$. Need to show that
    $0^3 + 11\cdot 0$ is divisible by $6$.
    \begin{compactenum}
    \item By arithmetics, we have $0^3 + 11\cdot 0 = 0$, which is
      divisible by $6$ ($0 = 6\cdot 0$).
    \end{compactenum}
  \item \textbf{Inductive Case}: $n = k + 1$. Need to show that
    $(k + 1)^3 + 11(k + 1)$ is divisible by $6$.\\
    \textit{Induction Hypothesis}: $k^3 + 11k$ is divisible by $6$.
    \begin{compactenum}
    \item By algebra, we have
      $(k + 1)^3 + 11(k + 1) = k^3 + 3k^2 + 3k + 1 + 11k + 11 =
      k^3 + 3k^2 + 14k + 12$.
    \item By the induction hypothesis there exists some $k'$ such
      that $k = 6k'$. Thus we have
      $k^3 + 3k^2 + 14k + 12 = (6k')^3 + 3(6k')^2 + 14(6k') + 12$.
    \item By algebra, we have
      $(6k')^3 + 3(6k')^2 + 14(6k') + 12 =
       6(6^2k'^3) + 6(3\cdot 6k'^2) + 6(14k') + 6\cdot 2 =
       6(6^2k'^3 + 3\cdot 6k'^2 + 14k' + 2)
      $, which is divisible by $6$.
    \end{compactenum}
    $\square$
  \end{itemize}

\item To avoid having to flip between the documents, here is the
  BNF grammar for lists and the definitions of $\mathit{length}$
  and $\mathit{concat}$
  \[
    l ::= \epsilon ~|~ x;l
  \]
  \[
  \mathit{length}(l) =
  \begin{cases}
    0 & \text{if } l = \epsilon\\
    1 + \mathit{length}(l')& \text{if } l = x; l'\\
  \end{cases}
  \]
  \[
  \mathit{concat}(l_1, l_2) =
  \begin{cases}
    l_2 & \text{if } l_1 = \epsilon\\
    x; \mathit{concat}(l_1', l_2) & \text{if } l_1 = x; l_1'\\
  \end{cases}
  \]
  We want to show that the length of the concatenation of two
  lists equals the sum of their individual lengths, or formally:
  \[
    \forall l_1,l_2. length(concat(l_1, l_2)) = length(l_1) + length(l_2)
  \]
  We start by assuming that we have some $l_1$ and $l_2$ and
  proceed to show
  $length(concat(l_1, l_2)) = length(l_1) + length(l_2)$ by
  induction on $l_1$:
  \begin{itemize}
  \item \textbf{Base Case}: $l_1 = \epsilon$. Need to show $\mathit{length}(\mathit{concat}(\epsilon, l_2)) = \mathit{length}(\epsilon) + \mathit{length}(l_2)$
    \begin{compactenum}
    \item By the definition of $\mathit{concat}$, we have $concat(\epsilon, l_2) = l_2$, and thus
      $\mathit{length}(\mathit{concat}(\epsilon, l_2)) = \mathit{length}(l_2)$.
    \item Since $\mathit{length}(\epsilon) = 0$, we have
      $\mathit{length}(\mathit{concat}(\epsilon, l_2)) =
      \mathit{length}(l_2) = \mathit{length}(\epsilon) +
      \mathit{length}(l_2)$.
    \end{compactenum}
  \item \textbf{Inductive Case}: $l_1 = x; l_1'$. Need to show
    $\mathit{length}(\mathit{concat}(x;l_1', l_2)) = \mathit{length}(x;l_1') + \mathit{length}(l_2)$\\
    \textit{Induction Hypothesis}: $\mathit{length}(\mathit{concat}(l_1', l_2)) = \mathit{length}(l_1') + \mathit{length}(l_2)$
    \begin{compactenum}
    \item By the definition of $\mathit{concat}$, we have
      we have $\mathit{concat}(x;l_1', l_2) = x;\mathit{concat}(l_1', l_2)$
      and thus $\mathit{length}(\mathit{concat}(x;l_1', l_2)) = \mathit{length}(x;\mathit{concat}(l_1', l_2))$.
    \item By the definition of $\mathit{length}$ we have $\mathit{length}(x;\mathit{concat}(l_1', l_2)) = 1 + \mathit{length}(\mathit{concat}(l_1', l_2))$.
    \item By the induction hypothesis we have
      $\mathit{length}(\mathit{concat}(l_1', l_2)) =
      \mathit{length}(l_1') + \mathit{length}(l_2)$.
    \item By the definition of $\mathit{length}$ we have $\mathit{length}(x;l_1') = 1 + \mathit{length}(l_1)$.
    \item Together, we have $\mathit{length}(\mathit{concat}(x;l_1', l_2)) = 1 + \mathit{length}(l_1') + \mathit{length}(l_2) = \mathit{length}(x;l_1') + \mathit{length}(l_2)$.
    \end{compactenum}
    $\square$
  \end{itemize}

\item We start by formalising that a filtered list only contains
  elements that satisfies the predicate:
  \[
    \forall x, l. \mathit{member}(x, \mathit{filter}(l, P)) \implies P(x)
  \]

  In order to prove this we start by assuming that we have some
  $x$ and $l$ such that
  $\mathit{member}(x, \mathit{filter}(l, P))$. We show $P(x)$ by
  induction over $l$:
  \begin{itemize}
  \item \textbf{Base Case}: $l = \epsilon$. Need to show $P(x)$
    \begin{compactenum}
    \item By assumption, we have
      $\mathit{member}(x, \mathit{filter}(\epsilon, P))$, but
      $\mathit{filter}(\epsilon, P) = \epsilon$ so we have
      $\mathit{member}(x, \epsilon)$. None of the derivation rules
      for $\mathit{member}$ allows an empty, list so we have a
      contradiction. The case holds vacuously.
    \end{compactenum}
  \item \textbf{Inductive Case}: $l = y; l'$. Need to show $P(x)$\\
    \textit{Induction Hypothesis}: $\mathit{member}(x, \mathit{filter}(l', P)) \implies P(x)$
    \begin{compactenum}
    \item There are two cases of $\mathit{filter}$ to consider.
      \begin{compactenum}
      \item If $P(y)$ holds, then $\mathit{filter}(y;l', P) = y';\mathit{filter}(l', P)$.
        By inversion of our assumption $\mathit{member}(x, y';\mathit{filter}(l, P))$
        we have either $x = y$, in which case $P(x)$ holds by the assumption in this case,
        or we have $\mathit{member}(x, \mathit{filter}(l', P))$, in which case $P(x)$ holds
        by the induction hypothesis.
      \item If $P(y)$ does not hold, then $\mathit{filter}(y;l', P) = \mathit{filter}(l', P)$,
        and $P(x)$ holds by the induction hypothesis with our assumption $\mathit{member}(x, \mathit{filter}(l', P))$.
      \end{compactenum}
    \end{compactenum}
    $\square$
  \end{itemize}

\item We want to show
  \[
    \forall l_1, l_2. \mathit{filter}(\mathit{concat}(l_1, l_2), P) = \mathit{concat}(\mathit{filter}(l_1, P), \mathit{filter}(l_2, P))
  \]
  We assume that we have some lists $l_1$ and $l_2$ and proceed by
  induction over $l_1$:
  \begin{itemize}
  \item \textbf{Base Case}: $l_1 = \epsilon$. Need to show $\mathit{filter}(\mathit{concat}(\epsilon, l_2), P) = \mathit{concat}(\mathit{filter}(\epsilon, P), \mathit{filter}(l_2, P))$
    \begin{compactenum}
    \item By the definition of $\mathit{concat}$, we have $concat(\epsilon, l_2) = l_2$, and thus
      $\mathit{filter}(\mathit{concat}(\epsilon, l_2), P) = \mathit{filter}(l_2, P)$.
    \item By the definition of $\mathit{filter}$, we have
      $\mathit{filter}(\epsilon, P)) = \epsilon$, and thus
      $\mathit{concat}(\mathit{filter}(\epsilon, P), \mathit{filter}(l_2, P)) =
      \mathit{concat}(\epsilon, \mathit{filter}(l_2, P))$
    \item By the definition of $\mathit{concat}$, we have
      $\mathit{concat}(\epsilon, \mathit{filter}(l_2, P)) =
      \mathit{filter}(l_2, P)$.
    \item Together, we have $\mathit{filter}(\mathit{concat}(\epsilon, l_2), P) = \mathit{filter}(l_2, P) = \mathit{concat}(\mathit{filter}(\epsilon, P), \mathit{filter}(l_2), P)$.
    \end{compactenum}
  \item \textbf{Inductive Case}: $l_1 = x; l_1'$. Need to show
    $\mathit{filter}(\mathit{concat}(x;l_1, l_2), P) = \mathit{concat}(\mathit{filter}(x;l_1, P), \mathit{filter}(l_2, P))$\\
    \textit{Induction Hypothesis}: $\mathit{filter}(\mathit{concat}(l_1, l_2), P) = \mathit{concat}(\mathit{filter}(l_1, P), \mathit{filter}(l_2, P))$
    \begin{compactenum}
    \item By the definition of $\mathit{concat}$, we have
      we have $\mathit{concat}(x;l_1', l_2) = x;\mathit{concat}(l_1', l_2)$
      and thus $\mathit{filter}(\mathit{concat}(x;l_1', l_2), P) = \mathit{filter}(x; \mathit{concat}(l_1', l_2), P)$.
    \item There are two cases of filter to consider:
      \begin{compactenum}
      \item If $P(x)$ holds, then $\mathit{filter}(x; \mathit{concat}(l_1', l_2), P) = x; \mathit{filter}(\mathit{concat}(l_1', l_2), P)$.
        By the induction hypothesis we have $x; \mathit{filter}(\mathit{concat}(l_1', l_2), P) =
        x; \mathit{concat}(\mathit{filter}(l_1, P), \mathit{filter}(l_2, P))$,
        which by the definitions of $\mathit{concat}$ and $\mathit{filter}$ is equal to
        $\mathit{concat}(\mathit{filter}(x;l_1, P), \mathit{filter}(l_2, P))$.
      \item If $P(x)$ does not hold, then $\mathit{filter}(x; \mathit{concat}(l_1', l_2), P) =
        \mathit{filter}(\mathit{concat}(l_1', l_2), P)$ which by the induction
        hypothesis is equal to $\mathit{concat}(\mathit{filter}(l_1, P), \mathit{filter}(l_2, P))$.
        By the definitions of $\mathit{concat}$ and $\mathit{filter}$, we have
        $\mathit{concat}(\mathit{filter}(l_1, P), \mathit{filter}(l_2, P)) =
         \mathit{concat}(\mathit{filter}(x;l_1, P), \mathit{filter}(l_2, P))$.
      \end{compactenum}
    \end{compactenum}
    $\square$
  \end{itemize}


\end{enumerate}

\end{document}
    \item By the induction hypothesis we have $\mathit{filter}(x; \mathit{concat}(l_1', l_2), P) = x;\mathit{concat}(\mathit{filter}(l_1, P), \mathit{filter}(l_2), P)$.
